% Bagian: second-order-logic
% Bab: syntax-and-semantics
% Subbagian: introduction

\documentclass[../../../include/open-logic-section]{subfiles}

\begin{document}

\olfileid[id]{sol}{syn}{int}

\olsection{Pengantar}

Dalam logika orde pertama, kita menggabungkan simbol-simbol nonlogis suatu
bahasa tertentu, yakni !!{constant}s, !!{function}s, dan !!{predicate}s yang
dimilikinya, dengan simbol-simbol logis untuk mengungkapkan hal-hal mengenai
!!{structure}s orde pertama. Hal ini dilakukan dengan menggunakan gagasan
keterpenuhan, yang menghubungkan !!a{structure}~$\Struct{M}$, bersama suatu
penugasan variabel~$s$, dengan !!a{formula}~$!A$: $\Sat{M}{!A}[s]$ berlaku jika
dan hanya jika apa yang diungkapkan oleh~$!A$ adalah benar ketika
!!{constant}s, !!{function}s, dan !!{predicate}s yang muncul di dalamnya
ditafsirkan sebagaimana ditentukan oleh~$\Struct{M}$, dan variabel-variabel
bebasnya ditafsirkan sebagaimana ditentukan oleh~$s$. Interpretasi
!!{identity}~$\eq$ sudah melekat
dalam definisi $\Sat{M}{!A}[s]$, demikian pula interpretasi~$\lforall$ dan
$\lexists$. Yang pertama selalu ditafsirkan sebagai relasi identitas pada
!!{domain}~$\Domain{M}$ struktur tersebut, dan kuantor selalu ditafsirkan
sebagai merentang atas seluruh !!{domain}. Namun, yang sangat penting,
kuantifikasi hanya diperbolehkan atas anggota-anggota !!{domain}, sehingga
hanya !!{variable}s objek yang diperbolehkan mengikuti suatu kuantor.

Dalam logika orde kedua, baik bahasa maupun definisi keterpenuhan diperluas
sehingga mencakup variabel fungsi dan predikat bebas maupun terikat, serta
kuantifikasi atas variabel-variabel tersebut. Variabel-variabel ini
berhubungan dengan !!{function}s dan !!{predicate}s dengan cara yang sama
seperti variabel objek berhubungan dengan !!{constant}s. Variabel-variabel
tersebut memainkan peran yang sama dalam pembentukan suku dan !!{formula}s
logika orde kedua, dan kuantifikasi atasnya ditangani dengan cara serupa.
Dalam semantik \emph{standar}, kuantor orde kedua merentang atas semua objek
yang mungkin dari tipe yang tepat (fungsi beraritas~$n$ dari $\Domain{M}^n$
ke~$\Domain{M}$ bagi variabel fungsi, dan relasi beraritas~$n$ bagi variabel
predikat). Misalnya, meskipun
$\lforall[\Obj{v_0}][(\Obj{P^1_0}(\Obj{v_0}) \lor \lnot
  \Obj{P^1_0}(\Obj{v_0}))]$ merupakan suatu formula baik dalam logika orde
pertama maupun dalam logika orde kedua, dalam logika orde kedua kita juga
dapat mempertimbangkan
$\lforall[\Obj{V^1_0}][\lforall[\Obj{v_0}][(\Obj{V^1_0}(\Obj{v_0}) \lor
    \lnot \Obj{V^1_0}(\Obj{v_0}))]]$ dan
$\lexists[\Obj{V^1_0}][\lforall[\Obj{v_0}][(\Obj{V^1_0}(\Obj{v_0}) \lor
    \lnot \Obj{V^1_0}(\Obj{v_0}))]]$. Karena kedua formula ini tidak memuat
variabel bebas, keduanya merupakan !!{sentence}s logika orde kedua. Di sini,
$\Obj{V^1_0}$ merupakan variabel predikat orde kedua beraritas~$1$.
Interpretasi yang diperbolehkan bagi~$\Obj{V^1_0}$ sama dengan yang dapat kita
tetapkan kepada suatu !!{predicate} beraritas~$1$ seperti~$\Obj{P^1_0}$, yakni
himpunan-himpunan bagian dari~$\Domain{M}$. Kuantifikasi atasnya lalu berarti
menyatakan bahwa
$\lforall[\Obj{v_0}][(\Obj{V^1_0}(\Obj v_0) \lor \lnot
  \Obj{V^1_0}(\Obj{v_0}))]$ berlaku bagi semua cara menetapkan suatu himpunan
bagian dari~$\Domain{M}$ sebagai nilai~$\Obj{V^1_0}$, atau bagi setidaknya satu
cara. Karena setiap himpunan memuat atau tidak memuat suatu objek tertentu,
kedua kalimat itu benar dalam setiap !!{structure}.
\end{document}
